% main.tex — your rrxiv paper.
%
% This template ships a one-claim "hello world" paper that compiles
% out-of-the-box. Edit the metadata at the top, then replace the body
% with your actual content. Keep using the `claim`, `evidence`,
% `observation`, `openquestion`, and `scope` environments — the rrxiv
% parser extracts these into the Canonical Intermediate Representation.
%
% Build:
%   ./scripts/build.sh
% or:
%   tectonic -X compile paper/main.tex --outdir build

\documentclass{rrxiv}

% --- rrxiv metadata (mirrored in rrxiv-meta.json) ---
\rrxivid{rrxiv-paper-template}        % leave as template id until first submission
\rrxivversion{v1}
\rrxivprotocolversion{0.1.0}
\rrxivlicense{CC-BY-4.0}
\rrxivtopics{example,template}

\title{A starter rrxiv paper}
\author{Your Name}
\date{\today}

\begin{document}
\maketitle

\begin{abstract}
This template is a working rrxiv submission. It contains one claim, one
piece of supporting evidence, one open question, and a citation, so it
exercises every layer of the protocol and validates against the schema
without further edits. Replace this abstract, the body, and the
\texttt{rrxiv-meta.json} metadata to start your own paper.
\end{abstract}

\section{Background}
\label{sec:background}

Replace this section with your motivation. Cite prior work via
\verb|\cite{key}| referring to entries in \texttt{paper/refs.bib}.

\section{Claims}
\label{sec:claims}

\begin{claim}[Template renders cleanly]
\label{claim:hello}
A fresh clone of \texttt{rrxiv-paper-template} compiles to a valid
rrxiv paper without modification: the PDF builds, the sidecar
\texttt{*.rrxiv.aux} emits, and the resulting CIR validates against
\texttt{cir.schema.json} v0.1.0~\cite{rrxiv-cir-schema}.
\end{claim}

\begin{evidence}[CI run on every push]
\label{ev:ci}
The workflow in \texttt{.github/workflows/build.yml} runs
\texttt{scripts/build.sh}, \texttt{scripts/extract-cir.sh}, and
\texttt{scripts/verify.sh} on every push and pull request. The badge in
\texttt{README.md} reflects the latest CI outcome.
\end{evidence}

\begin{openquestion}[How should template versions be tracked?]
\label{oq:versioning}
This template is intentionally simple; future revisions may add
recommended figure conventions, a \texttt{paper/appendix.tex} include
pattern, or a script for vendored-class refresh. The right cadence and
governance for those updates is an open process question for the rrxiv
project.
\end{openquestion}

\bibliographystyle{plainnat}
\bibliography{refs}

\end{document}
